\par\noindent\textbf{(1)} 自由演化纯态相位为(|0>+e\textasciicircum{}\{iφ\}|1>)/√2；纯退相干通道使ρ(t)=(1/2)[[1,e\textasciicircum{}\{−Γφt\}e\textasciicircum{}\{−iφ\}],[e\textasciicircum{}\{−Γφt\}e\textasciicircum{}\{iφ\},1]]，迹保持且本征值(1±e\textasciicircum{}\{−Γφt\})/2非负。Bloch矢量长度e\textasciicircum{}\{−Γφt\}逐步收缩。
\par\noindent\textbf{(2)} 测量旋转后的赤道投影，选择分析相位ϑ时二值事件概率p\_+=(1+V cos(φ−ϑ))/2，其中可见度V=e\textasciicircum{}\{−Γφt\}。另一事件p\_−=1−p\_+，相位信息主要来自干涉曲线斜率大的象限而不是峰顶。
\par\noindent\textbf{(3)} 二值Fisher信息F\_φ=(∂φp\_+)²/[p\_+(1−p\_+)]=V² sin²(φ−ϑ)/[1−V² cos²(φ−ϑ)]。在ϑ=φ±π/2最佳，F\_φ,max=V²=e\textasciicircum{}\{−2Γφt\}；退相干损失信息而非增加相位灵敏度。
\par\noindent\textbf{(4)} 若估计频率失谐Δ，因φ=Δt，单次F\_Δ=t²e\textasciicircum{}\{−2Γφt\}，在固定每次实验、忽略制备开销时对t求极值得t*=1/Γφ；若固定总实验时长，实验次数约T\_total/t使总信息∝T\_total t e\textasciicircum{}\{−2Γφt\}，最优改为t=1/(2Γφ)，必须说明资源约束。
